% peano-2606-bdb
% https://docs.google.com/document/d/1fP0LCSwN_x0zG9ieAL4_rlilmyP3mgOxQ-HRaDrEJUI/edit?tab=t.0
% !TEX encoding = UTF-8 Unicode
% https://github.com/matheusgirola/Halmos-Naive-Set-Theory-OCR-LaTeX-Reedition/tree/master
% bdb: See https://gemini.google.com/app/ae3a86e901b911fb. bda: Even more step by step. Steps: order, total, well. bcd: will change Lemma {Ltot}. bcc: will try to go more step by step. bcb: Will add a prop. stating the trichotomy. bca: Back to proving antisymmetry first (sigh). bbc: Planning to derive antisymmetry of \le from trichotomy. Versions peano-2606-bba and peano-2606-bbb were ok.
\documentclass[12pt,letterpaper]{article}
\usepackage[T1]{fontenc}
\usepackage{amssymb,amsmath,amsthm} 
\usepackage[letterpaper,top=40pt,left=50pt,right=50pt,bottom=60pt]{geometry} % print
% scribe https://docs.google.com/document/d/1DJcZrjPx_dnQfeFpTTFppo3rdrFN6FbktpK75whE2zQ/edit?tab=t.0
%\usepackage[letterpaper,total={109mm,145mm},centering]{geometry} % SMALL4 scribe 
%\usepackage[letterpaper,total={101mm,135mm},centering]{geometry} % SMALL5 scribe last used
\usepackage{microtype}
%\usepackage{tikz-cd}\usepackage{varwidth}
%\usepackage{comment}
\usepackage[pdfusetitle]{hyperref}
%\usepackage{marvosym}% emoji https://tex.stackexchange.com/questions/3695/smileys-in-latex https://ctan.org/pkg/marvosym?lang=en
\usepackage{biblatex} 
\addbibresource{halmos-2606.bib} 
%\pagestyle{empty}
%\setlength\parindent{0pt}
%\setlength{\parskip}{3pt} % variable
%\renewcommand{\baselinestretch}{1.1} % variable
%\newcommand{\noi}{\noindent}
\newtheorem{thm}{Theorem}%[section] \newtheorem{cor}[thm]{Corollary}
\newtheorem{lem}[thm]{Lemma} %\newtheorem{remark}[thm]{Remark}
\newtheorem{prop}[thm]{Proposition}
\title{Peano systems} \author{Pierre-Yves Gaillard} \date{}

\begin{document}% \tiny \scriptsize \footnotesize \small \normalsize \large \Large \LARGE \huge \Huge \Gamma \Delta \Theta \Lambda \Xi \Pi \Sigma \Phi \Psi \Omega

\maketitle %\tableofcontents
%\centerline{\scriptsize This text is available at \url{https://github.com/Pierre-Yves-Gaillard/peano-systems}.} 
\begin{center}
{\scriptsize This text is available at \url{https://github.com/Pierre-Yves-Gaillard/peano-systems}.}
\end{center}%\bigskip 

The purpose of this text is to rewrite some parts of Halmos's and Enderton's set theory books \cite{halmos_naive_1974} and \cite{Enderton1977}. More precisely we rewrite parts of Chapters~12 and 13 of \cite{halmos_naive_1974} and parts of Chapter~4 of \cite{Enderton1977}. This means in particular that, in this text, \emph{we only use what has been defined and proved in the previous chapters of each book}. (For instance, the addition of natural numbers is taken for granted in usual mathematics, but not here.) At this point of the two books, the set \(\omega\) of natural numbers has just been defined (as a mere set) by virtue of the infinity axiom, and almost none of the familiar properties of \(\omega\) has been established. The word ``rewrite'' should be taken in a broad sense: we try not only to prove these properties, but also to put them in a more general setting. 

We use the following definitions: A \textbf{system} is a triple \(X^*=(X,\xi,x_0)\) where \(X\) is a set, \(\xi\) is a map from \(X\) to itself and \(x_0\) is an element of \(X\). We sometimes abuse the terminology and say the set \(X\) itself is a system when \(\xi\) and \(x_0\) are obvious from the context. A subset \(Y\subset X\) is \(\xi\)-\textbf{closed} if \(\xi(Y)\subset Y\) (here and in the sequel \(U\subset V\) means ``\(U\) is a subset of \(V\)''). If \(x\in X\) then the intersection of all \(\xi\)-closed subsets of \(X\) containing \(x\) is a \(\xi\)-closed subset denoted by \(X_x\) and is called the \(\xi\)-closed subset \textbf{generated} by \(x_0\); it is the least \(\xi\)-closed subset containing \(x\). A \textbf{morphism} from the system \(X^*=(X,\xi,x_0)\) to the system \(Y^*=(Y,\eta,y_0))\) is a map \(f:X\to Y\) such that \(f(x_0)=y_0\) and \(f(\xi(x))=\eta(f(x))\) for all \(x\in X\). We say that \(f\) is an \textbf{isomorphism} if it is bijective; in this case, \(f^{-1}\) is a morphism from \(Y^*\) to \(X^*\), and we write \(X^*\simeq Y^*\). The identity of \(X\) is a morphism from \(X^*\) to itself. If \(g:Y\to Z\) is a morphism from \(Y^*\) to \(Z^*=(Z,\zeta,z_0))\), then \(g\circ f\) is a morphism from \(X^*\) to \(Z^*\). 

A system \(A^*=(A,\alpha,a_0)\) is \textbf{initial} if for all system \(X^*=(X,\xi,x_0)\) there is exactly one morphism from \(A^*\) to \(X^*\). We say that the system \(X^*\) is \textbf{Peano} if 
\begin{equation}\label{Edisj}
x_0\notin \xi(X),
\end{equation}
\begin{equation}\label{Einj}
\xi\text{ is injective},
\end{equation}
\begin{equation}\label{Eind}
X_{x_0}=X. 
\end{equation} 
(Recall that \(X_{x_0}\) is the \(\xi\)-closed subset generated by \(x\).) Property~\eqref{Eind} is known as the \textbf{principle of mathematical induction}. The terminology ``Peano system'' is used by Enderton in \cite{Enderton1977}. 

The following obvious fact will be implicitly used throughout: Let \(X^*\) and \(Y^*\) be isomorphic. If \(X^*\) is initial, so is \(Y^*\); if \(X^*\) is Peano, so is \(Y^*\). 

At the beginning of Chapter 12 in \cite{halmos_naive_1974} and of Chapter 4 in \cite{Enderton1977} a system \(\omega^*=(\omega,n\mapsto n^+,0)\) has been defined, where \(\omega,n\mapsto n^+\) and \(0\) are intended to represent respectively the set of natural numbers, the successor map \(n\mapsto n+1\), and the number \(0\). We want to prove: 

\begin{thm}\label{T1} 
Consider the system \(\omega^*=(\omega,n\mapsto n^+,0)\) and arbitrary systems \(X^*=(X,\xi,x_0)\) and \(Y^*=(Y,\eta,y_0)\). Then we have: 
\begin{itemize}
\item[\emph{(a)}] \(\omega^*\) is Peano (in particular Peano systems exist), 
\item[\emph{(b)}] \(X^*\) is Peano if and only if it is initial, 
\item[\emph{(c)}] if \(X^*\) is Peano, then there is a unique well-order \(\le\) on \(X\) such that \(x\le y\Leftrightarrow y\in X_x\Leftrightarrow X_y\subset X_x\) for all \(x,y\in X\),
\item[\emph{(d)}] if \(X^*\) and \(Y^*\) are initial, then the unique morphism from \(X^*\) to \(Y^*\) is an isomorphism.   
\end{itemize} 
(Recall that \(X_x\) is the \(\xi\)-closed subset of \(X\) generated by \(x\), and that a \textbf{well-order} is an order \(\le\) on a set \(Z\) such that any nonempty subset of \(Z\) has a least element.) 
\end{thm} 

Part (a) is proved by Halmos and Enderton in \cite{halmos_naive_1974} and \cite{Enderton1977} respectively. 

\begin{lem}\label{Ld} 
Part \emph{(d)} in Theorem~\ref{T1} holds. That is, if \(X^*\) and \(Y^*\) are initial, then the unique morphism from \(X^*\) to \(Y^*\) is an isomorphism. 
\end{lem} 
\begin{proof}  
Let \(f: X^* \to Y^*\) and \(g: Y^* \to X^*\) be the unique morphisms guaranteed by initiality. The composition \(g \circ f: X^* \to X^*\) is a morphism from \(X^*\) to itself. Since the identity map \(\text{id}_X\) is also a morphism from \(X^*\) to itself, uniqueness dictates that \(g \circ f = \text{id}_X\). By a symmetric argument, \(f \circ g = \text{id}_Y\), proving that \(f\) is an isomorphism.
\end{proof} 

\section{\texorpdfstring{\(\omega^*\)}{ω*} is Peano} %%%%%%%%%

\begin{prop}\label{Pp} 
The system \(\omega^*=(\omega,n\mapsto n^+,0)\) is Peano. 
\end{prop} 

Recall that we admit only what has been defined and proved up to the beginning of Chapter~12 of \cite{halmos_naive_1974} and of Chapter~4 of \cite{Enderton1977}, and that, at this point of the two books, we have defined the \textbf{successor} of an arbitrary set \(x\) as the set \(x^+=x\cup\{x\}\). We have also defined a set \(\omega\) such that \(\varnothing\in\omega\), where \(\varnothing\) is the empty set. When \(\varnothing\) is regarded as an element of \(\omega\), it is usually denoted by 0. We have shown \(n^+\in\omega\) whenever \(n\in\omega\). Let us denote the successor map \(\omega\to\omega\) by \(n\mapsto n^+\) or by \(s\), whatever is more convenient. 

We also know that, in the above notation, we have \(\omega_0=\omega\), that is, the principle of mathematical induction holds. Since condition~\eqref{Edisj} above, namely \(0\notin s(\omega)\), is obvious, it only remains to verify the injectivity of \(s\) in order to prove that \(\omega^*=(\omega,n\mapsto n^+,0)\) is Peano. To check this injectivity, we first note that, for all \(i,j\in\omega\), we have 
\begin{equation}\label{Emain}
j\in j^+,\ j\subset j^+\text{ and }i\in j^+\Leftrightarrow (i\in j\text{ or }i=j). 
\end{equation} 

\begin{lem}\label{Lmain} 
For all \(i,j,k\in\omega\) we have \(i\in j\in k\Rightarrow i\in k\) or, equivalently, for all \(m,n\in\omega\) we have \(m\in n\Rightarrow m\subset n\). 
\end{lem} 
\begin{proof} 
Given \(i,j\in\omega\) we denote by \(\Gamma(k)\) the statement \(i\in j\in k\Rightarrow i\in k\), and we prove \(\Gamma(k)\) for all \(k\) by induction. Note that \(\Gamma(0)\) is true, because it is of the form ``False implies something''. To prove \(\Gamma(k)\Rightarrow \Gamma(k^+)\), we assume \(\Gamma(k)\) and \(i\in j\in k^+\). We must show \(i\in k^+\). By \eqref{Emain}, we have \(j\in k\) or \(j=k\), and we want to prove \(i\in k\) or \(i=k\). We easily see that we have in fact \(i\in k\) in both cases. 
\end{proof} 

\begin{proof}[Proof of Proposition~\ref{Pp}] 
As already observed, it only remains to verify the injectivity of \(n\mapsto n^+\). Assuming \(i\ne j\) and \(i^+=j^+\) for the sake of contradiction, we get \(i\in j\in i\) by \eqref{Emain}, hence \(i\subset j\subset i\) by Lemma~\ref{Lmain} (implication \(m\in n\Rightarrow m\subset n\)), hence \(i=j\). 
\end{proof} 

\section{Well-order} %%%%%%%%% 

We want to prove Part (c) of Theorem~\ref{T1}, which we state below (with a slightly different notation) as Proposition~\ref{Pc}. 

\begin{prop}\label{Pc} 
If \(U^*=(U,s,0)\) is a Peano system, then there is a unique well-order \(\le\) on \(U\) such that \(u\le v\Leftrightarrow U_v\subset U_u\) for all \(u,v\in U\). 
\end{prop} 

Let \(U^*=(U,s,0)\) be a Peano system. Set \(U'=U\setminus\{0\}\). 

\begin{lem}\label{Lbij} 
We have \(s(U)=U'\), that is, \(s\) induces a bijection \(s':U\to U'\). 
\end{lem} 
\begin{proof} 
We have \(0\notin s(U)\) by \eqref{Edisj} and \(U=\{0\}\cup s(U)\) by induction (see \eqref{Eind}). 
\end{proof} 

The following lemma, which is an obvious consequence of Lemma~\ref{Lbij}, will be useful: 

\begin{lem}\label{Liso} 
The bijection \(s':U\to U'\) of Lemma~\ref{Lbij} induces an isomorphism of systems from \(U^*=(U,s,0)\) to \(U'^*=(U',s'',s(0))\), where \(s'':U'\to U'\) is induced by \(s\). (In particular \(U_{s(0)}=U'=s(U)\), \(U'^*\) is Peano and \(U_{s(u)}=s(U_u)\) for all \(u\).) 
\end{lem} 

In the sequel we abuse the notation and designate \(s'\) and \(s''\) simply by \(s\) whenever convenient. (Note that \(s^{-1}(u)\in U\) is well-defined for \(u\in U'\)---but not for \(u=0\).) 

\begin{lem}\label{Lgen} 
Let \(X\) be a subset of the product \(U\times U\) containing \(U\times\{0\}\) and \(\{0\}\times U\) (as subsets), and containing \((s(u),s(v))\) whenever it contains \((u,v)\) (as elements). Then \(X=U\times U\). 
\end{lem} 
\begin{proof} 
Let \(Y\) be the set of all \(u\in U\) such that \((u,v)\in X\) for all \(v\in U\). Then \(0\) is in \(Y\) because \(X\) contains \(\{0\}\times U\). Let \(u\) be in \(Y\) and \(v\) in \(U\). It suffices to show \((s(u),v)\in X\). By Lemma~\ref{Lbij} we have \(v=0\) or \(v=s(w)\) for some \(w\). The conclusion follows in the first case from the fact that \(X\) contains \(U\times\{0\}\), and in the second case from the fact that \(X\) contains \((s(u),s(w))\) because it contains \((u,w)\) (because \(u\in Y\)). 
\end{proof} 

We define the relation \(\le\) on \(U\) by setting \(u\le v\Leftrightarrow U_v\subset U_u\), or, equivalently, \(u\le v\Leftrightarrow v\in U_u\). 

\begin{lem}\label{Lpo} 
 \begin{itemize} 
  \item[\emph{(a)}] The relation \(\le\) is a partial order on \(U\). 
  \item[\emph{(b)}] The bijection \(s':U\to U'\) of Lemma~\ref{Lbij} is an isomorphism of partially ordered sets (here \(U'\) is endowed with the subset order). 
  \item[\emph{(c)}] For all \(u,v\in U\) we have \(U_u=U_v\Rightarrow u=v\). 
 \end{itemize} 
\end{lem} 
\begin{proof} 
It suffices to prove (c). For all \(u,v\in U\) let \(\Delta(u,v)\) be the statement \(U_u=U_v\Rightarrow u=v\). Note that \(\Delta(u,v)\) is symmetric in \(u\) and \(v\), that the implication \(\Delta(u,v)\Rightarrow\Delta(s(u),s(v))\) follows immediately from Lemma~\ref{Liso}, and that \(\Delta(u,u)\) is obvious. Thus, it suffices, by Lemmas~\ref{Lbij} and \ref{Lgen}, to prove \(\Delta(s(u),0)\), that is, for all \(u\) we have \(U_{s(u)}=U\Rightarrow s(u)=0\). But we get \(U_{s(u)}=s(U_u)\subset U'\) by Lemma~\ref{Liso}, hence \(U_{s(u)}\ne U\), hence we only need to prove that ``False implies \(s(u)=0\)'', which is true (because ``False'' implies anything). 
\end{proof} 

\begin{lem}\label{Ltot} 
The set \(U\) is \emph{totally} ordered by \(\le\), that is, for all \(u,v\in U\), we have \(u\le v\) or \(v\le u\). 
\end{lem} 
\begin{proof} 
Let \(\Theta(u,v)\) be the statement: \(u\le v\) or \(v\le u\). Since \(\Theta(u,v)\) is symmetric in \(u\) and \(v\), Lemmas~\ref{Lgen} and \ref{Lpo}(b) imply that it suffices to prove \(\Theta(u,0)\). But this follows from the fact that we have 
\begin{equation}\label{E0min}
u\ge0
\end{equation} 
for all \(u\in U\). Indeed \eqref{E0min} is equivalent to \(u\in U_0\) by definition, and the equality \(U_0=U\) follows from \eqref{Eind}. 
\end{proof} 

We shall often use Lemma~\ref{Ltot} in an implicit way. Given \(u\in U\), let \([0,u]\) be the set of all \(v\in U\) such that \(v\le u\) (the notation is justified by \eqref{E0min}). 

\begin{lem}\label{Lprelim} 
For all \(u,v\in U\) we have 
 \begin{itemize} 
  \item[\emph{(a)}] \(s(u)>u\),  
  \item[\emph{(b)}] \(v>u\Rightarrow v\ge s(u)\), 
  \item[\emph{(c)}] \([0,0]=\{0\}\),  
  \item[\emph{(d)}] \([0,s(u)]=[0,u]\cup\{s(u)\}\). 
 \end{itemize}
\end{lem} 
\begin{proof} 
(a) Let \(u\) be in \(U\). The subset \(U_u\) of \(U\) is \(s\)-closed, hence \(s(u)\in U_u\), hence \(s(u)\ge u\). We prove \(s(u)\ne u\) by induction. The case \(u=0\) and the induction step follow respectively from \eqref{Edisj} and \eqref{Einj}. Now (a) is proved. 

(b) Let \(\Lambda(u,v)\) be the statement \(v>u\Rightarrow v\ge s(u)\). By Lemma~\ref{Lgen} it is enough to verify that for all \(u,v\in U\) we have: (i) \(\Lambda(u,0)\), (ii) \(\Lambda(0,v)\) and (iii) \(\Lambda(u,v)\Rightarrow\Lambda(s(u),s(v))\). By \eqref{E0min}, the statement (i) is of the form ``False implies something'', which is true. To prove (ii), that is \(v>0\Rightarrow v\ge s(0)\), note that \(0<v<s(0)\) would yield \(s^{-1}(v)<0\) by Lemma~\ref{Lpo}(b), contradicting \eqref{E0min}. Finally, (iii) follows from Lemma~\ref{Lpo}(b). 

(c) This follows from \eqref{E0min}. 

(d) This follows from (a), (b) and (c). 
\end{proof} 

\begin{proof}[Proof of Proposition~\ref{Pc}] 
Note that an element of a totally ordered set is the least element if and only if it is minimal. Recall that \(U\) is totally ordered (Lemma~\ref{Ltot}). Assume that \(X\subset U\) has no minimal element. It is enough to show \(X=\varnothing\). It suffices to prove \([0,u]\cap X=\varnothing\) for all \(u\in U\). This is true for \(u=0\) by Lemma~\ref{Lprelim}(c). Assume, to the contrary, that it is true for \(u\) and false for \(s(u)\). We get \(s(u)\in X\) by Lemma~\ref{Lprelim}(d), hence, since \(s(u)\) is not minimal in \(X\), there is a \(v\in X\) which is less than \(s(u)\), hence \(v\in[0,u]\cap X\) again by Lemma~\ref{Lprelim}(d), contradiction. 
\end{proof} 

\section{Peano implies initial} %%%%%%%%% 

\begin{prop}\label{PP2i} 
Peano systems are initial. 
\end{prop} 

As above, we let \(U^*=(U,s,0)\) be a Peano system, and we will show that it is initial. Let \(X^*=(X,\xi,x_0)\) be an arbitrary system. Proposition~\ref{PP2i} will follow from the lemma below: 

\begin{lem}\label{LU2X} 
There is a unique map \(f:U\to X\) such that \(f(0)=x_0\) and \(f(s(u))=\xi(f(u))\) for all \(u\in U\).  
\end{lem} 

We prove the uniqueness in Lemma~\ref{LU2X}. Let \(g:U\to X\) satisfy the same conditions as \(f\), and assume, to the contrary, that \(g\ne f\). By Proposition~\ref{Pc}, Lemma~\ref{Lbij} and the equalities \(g(0)=x_0=f(0)\), there is a least \(u\in U\) such that \(g(s(u))\ne f(s(u))\). The minimality of \(u\) and the inequality \(s(v)>v\) (Lemma~\ref{Lprelim}(a)) imply \(g(u)=f(u)\), and thus \(g(s(u))=\xi(g(u))=\xi(f(u))=f(s(u))\), contradiction. 

\begin{lem}\label{Lrec} 
If \(X^*=(X,\xi,x_0)\) is an arbitrary system, then for all \(u\in U\) there is a unique map \(f_u:[0,u]\to X\) such that \(f_u(0)=x_0\), and \(f_u(s(v))=\xi(f_u(v))\) whenever \(s(v)\le u\). Moreover, the family \((f_u)_{u\in U}\) satisfies \(f_{s(u)}(s(u))=\xi(f_u(u))\) for all \(u\). 
\end{lem} 
\begin{proof} 
We first prove that for all \(u\) there is at most one such map \(f_u\). Let \(g:[0,u]\to X\) satisfy the same conditions as \(f_u\), and assume, to the contrary, that \(g\ne f_u\). By Proposition~\ref{Pc}, Lemma~\ref{Lbij} and the equalities \(g(0)=x_0=f_u(0)\), there is a least \(v\in U\) such that \(s(v)\le u\) and \(g(s(v))\ne f_u(s(v))\). The minimality of \(u\) and the inequality \(s(v)>v\) (Lemma~\ref{Lprelim}(a)) imply \(g(v)=f_u(v)\), and thus \(g(s(v))=\xi(g(v))=\xi(f_u(v))=f_u(s(v))\), contradiction. Now we prove the existence of such a \(f_u\) by induction on \(u\). We define first \(f_0\) by \(f_0(0)=x_0\) and then \(f_{s(u)}:[0,s(u)]\to X\) by extending \(f_u:[0,u]\to X\) (which exists and has the relevant properties by the induction hypothesis) to a map \([0,s(u)]=[0,u]\cup\{s(u)\}\to X\) (see Lemma~\ref{Lprelim}(d)) via the formula \(f_{s(u)}(s(u))=\xi(f_u(u))\). The equalities \(f_{s(u)}(0)=x_0\) and \(f_{s(u)}(s(v))=\xi(f_{s(u)}(v))\) for \(v\le u\) follow immediately from the definition of \(f_{s(u)}\) and the conditions that \(f_u\) satisfies under the induction hypothesis. 
\end{proof} 

We prove Lemma~\ref{LU2X}, and thus Proposition~\ref{PP2i}. It only remains to prove the existence of \(f\). Let \((f_u)_{u\in U}\) be the family given by Lemma~\ref{Lrec}. Then we define \(f\) by \(f(u)=f_u(u)\), and the theorem follows from Lemma~\ref{Lrec}. 

\section{Proof of Theorem~\ref{T1}} %%%%%%%%%%%%%%%

Recall the statements in Theorem~\ref{T1}: 

\begin{itemize}
\item[(a)] \(\omega^*\) is Peano, 
\item[(b)] \(X^*\) is Peano if and only if it is initial, 
\item[(c)] if \(X^*\) is Peano, then there is a unique well-order \(\le\) on \(X\) such that \(x\le y\Leftrightarrow y\in X_x\Leftrightarrow X_y\subset X_x\) for all \(x,y\in X\),
\item[(d)] if \(X^*\) and \(Y^*\) are initial, then the unique morphism from \(X^*\) to \(Y^*\) is an isomorphism.   
\end{itemize} 

Parts~(a), (c) and (d) have been respectively proved as Proposition~\ref{Pp}, Proposition~\ref{Pc} and Lemma~\ref{Ld}.  

\begin{proof}[Proof of \emph{(b)}]  
Peano systems are initial by Proposition~\ref{PP2i}. Let \(X^*\) be initial. By (d) and the fact that \(\omega^*\) is initial, \(X^*\) is isomorphic to \(\omega^*\). Then \(X^*\) is Peano because \(\omega^*\) is Peano by (a). 
\end{proof} 

\section{Finite sets} %%%%%%%%%%%%%%%

If \(X\) and \(Y\) are two sets, we indicate the existence of a bijection \(X\to Y\) by \(X\sim Y\), and we say that \(X\) and \(Y\) are \textbf{equipotent}. 

\begin{lem}\label{LabX} 
Let \(X\) be a set and let \(a,b\in X\). Then \(X\setminus\{a\}\sim X\setminus\{b\}\).
\end{lem} 
\begin{proof} 
We may assume \(a\ne b\). Then the map \(f:X\setminus\{a\}\to X\setminus\{b\}\) defined by \(f(x)=a\) if \(x=b\) and \(f(x)=x\) otherwise is bijective. 
\end{proof} 

\begin{lem}\label{LabXY} 
Let \(X\) and \(Y\) be two sets, and let \(a\in X,b\in Y\). Then \(X\setminus\{a\}\sim Y\setminus\{b\}\) if and only if \(X\sim Y\). In particular, for \(n\in\omega\) we have \(X\setminus\{a\}\sim n\) if and only if \(X\sim n^+\).
\end{lem} 
\begin{proof} 
We can extend any bijection \(X\setminus\{a\}\to Y\setminus\{b\}\) to a bijection \(X\to Y\) by mapping \(a\) to \(b\). Conversely, if \(f:X\to Y\) is bijective, then \(f(X\setminus\{a\})=Y\setminus\{f(a)\}\), and the result follows from Lemma~\ref{LabX}. 
\end{proof} 

\begin{lem}\label{Lbasic} 
For all \(i,j\in\omega\) we have \(i\sim j\Rightarrow i=j\). 
\end{lem} 
\begin{proof} 
Let \(\Xi(i)\) be the following statement: for all \(j\in\omega\) we have \(i\sim j\Rightarrow i=j\). Then \(\Xi(0)\) is clear. Let us prove \(\Xi(i^+)\) assuming \(\Xi(i)\). We have \(i^+\sim j\) and we must show \(i^+=j\). Since \(j\ne0\), Lemma~\ref{Lbij} implies \(j=k^+\) for some \(k\). Since \(i^+\sim k^+\), Lemma~\ref{LabXY} entails \(i\sim k\), and \(\Xi(i)\) yields \(i=k\), hence \(i^+=k^+=j\). 
\end{proof} 

We say that a set \(X\) is \textbf{finite} if it is equipotent to some \(n\in\omega\). Lemma~\ref{Lbasic} implies that a finite set is equipotent to a unique \(n\in\omega\). If this is the case, we write \(n=|X|\) and we say that \(n\) is the \textbf{number of elements} of \(X\). 

\begin{lem}\label{Lsubset} 
A subset of a finite set is finite. 
\end{lem} 
\begin{proof} 
It suffices to show that a subset \(X\) of \(n\in\omega\) is finite. We induct on \(n\). If \(n=0\), this is obvious. Assume that the statement holds for \(n\) and let us prove it for \(n^+\). So let \(X\subset n^+=n\cup\{n\}\). We may assume \(X\ne n^+\). Then \(X\) is contained in \(n^+\setminus\{i\}\) for some \(i\in n^+\), hence Lemma~\ref{LabXY} implies that \(X\) is equipotent to some subset \(Y\) of \(n=n^+\setminus\{n\}\), hence \(Y\) is finite by the induction hypothesis. 
\end{proof} 

\begin{lem}\label{Lunion} 
The union of two finite sets is finite. 
\end{lem} 
\begin{proof} 
We first consider the particular case when our finite sets are \emph{disjoint}. Given \(j\in\omega\), let \(\Pi(i)\) be the following statement: the conditions \(i\in\omega,i\sim X,j\sim Y\) and \(X\cap Y=\varnothing\) imply that \(X\cup Y\) is finite. Then \(\Pi(0)\) is obvious. Let us prove \(\Pi(i^+)\) assuming \(\Pi(i)\). We have \(i^+\sim X,j\sim Y,X\cap Y=\varnothing\), and we must show that \(X\cup Y\) is finite. Let \(a\) be in \(X\). Then \(X\setminus\{a\}\sim i\) by Lemma~\ref{LabXY}; hence, \((X\cup Y)\setminus\{a\}=(X\setminus\{a\})\cup Y\) is finite by \(\Pi(i)\), and thus \(X\cup Y\) is finite by Lemma~\ref{LabXY}. This proves the result when the sets are disjoint. If \(X\) and \(Y\) are arbitrary finite sets, \(X \cup Y\) can be written as the union of three pairwise disjoint sets: \(X \setminus Y\), \(X \cap Y\), and \(Y\setminus X\). Because each of these is a subset of either \(X\) or \(Y\), they are finite by Lemma~\ref{Lsubset}. The general case then follows directly from the disjoint case. 
\end{proof} 

Let \(i,j\in\omega\). It is easy to see that there are sets \(X\) and \(Y\) such that \(i\sim X,j\sim Y, X\cap Y=\varnothing\) (for example, \(X=i,Y=j\times 0^+\)), and that the element \(|X\cup Y|\in\omega\) (which exists by Lemma~\ref{Lunion}) depends only on \(i\) and \(j\), and not on the choice of \(X\) and \(Y\), so that we can define the \textbf{sum} \(i+j\in\omega\) of \(i\) and \(j\) by \(i+j=|X\cup Y|\). We call this operation \textbf{addition}. Then, properties such as associativity and commutativity of addition follow swiftly from the corresponding properties of the union of sets (which follow in turn from the corresponding properties of the Boolean operator \(\lor\)). One can also define multiplication and exponentiation along the same lines. 

\printbibliography

\end{document}
